


% \begin{table*}[t]
% \centering
% \small
% \caption{Per-instance results of BP-MDS ($\rho = 10^4$).}
% \label{tab:bpmds_per_instance_results}
% \begin{tabular}{|c|c|c|c|c|c|c|c|c|}
% \hline
% Class & N & Input & BKS & $\alpha$ ($^\circ$) & Time(s) & Memory(MB) & Cost & Gap(\%)\\
% \hline
% \noalign{\vskip 2pt} 
% \multirow{3}{*}{S} 
%  & 151 & CMT4 & 1028.42 & 127.0 & 0.0395 & 9.113 & 1143.98 & 11.23 \\
%  & 241 & Golden\_17 & 707.76 & 36.0 & 0.07 & 10.52 & 762.71 & 7.76 \\
%  & 301 & Golden\_18 & 995.13 & 24.0 & 0.12 & 14.88 & 1044.63 & 4.97 \\
% \hline

% \noalign{\vskip 2pt}
% \multirow{5}{*}{M} 
%  & 3001 & Leuven1 & 192848 & 30.0 & 0.302 & 10.26 & 207573.35 & 7.64 \\
%  & 4001 & Leuven2 & 111391 & 66.0 & 3.0 & 10.26 & 127447.35 & 14.41 \\
%  & 6001 & Antwerp1 & 477277 & 18.0 & 3.28 & 51.85 & 510570.898 & 6.97 \\
%  & 7001 & Antwerp2 & 291350 & 39.0 & 1.597 & 9.76 & 327747.64 & 12.49 \\
%  & 10001 & Ghent1 & 469531 & 23.0 & 0.78 & 10.41 & 503565.34 & 7.25 \\
% \hline

% \noalign{\vskip 2pt}
% \multirow{7}{*}{L} 
%  & 11001 & Ghent2 & 257749 & 35.5 & 2.04 & 10.26 & 292125.24 & 13.34 \\
%  & 15001 & Brussels1 & 501719 & 11.5 & 1.01 & 228.58 & 544650.63 & 8.56 \\
%  & 16001 & Brussels2 & 345481 & 31.0 & 2.25 & 9.75 & 390855.84 & 11.61 \\
%  & 20001 & Flanders1 & 7240124 & 18.0 & 1.34 & 11.47 & 7668875.48 & 5.92 \\
%  & 30001 & Flanders2 & 4373320 & 19.5 & 4.65 & 12.38 & 4864684.64 & 11.24 \\
%  & 25001 & XML25000\_1173\_01 & -- & 2.0 & 2.31 & 346.74 & 46902788.87 & -- \\
%  & 50001 & XML50000\_2173\_01 & -- & 2.5 & 5.04 & 343.13 & 58108045.17 & -- \\
% \hline

% \noalign{\vskip 2pt}
% \multirow{2}{*}{XL} 
%  & 75001 & XML75000\_1173\_01 & -- & 1.0 & 5.87 & 331.14 & 141858875.48 & \textunderscore{} \\
%  & 100001 & XML100000\_2173\_01 & -- & 1.0 & 9.73 & 333.01 & 116740181.51 & \textunderscore{} \\
% \hline

% \noalign{\vskip 2pt}
% \multirow{2}{*}{XXL}
%  & 500001 & XML500000\_1173\_01 & -- & 0.3 & 40.98 & 315.07 & 946129677.99 & \textunderscore{} \\
%  & 1000001 & XML1000000\_2123\_01 & -- & 0.3 & 88.67 & 288.31 & 993950684.86 & \textunderscore{} \\
% \hline

% \noalign{\vskip 2pt}
% \multirow{4}{*}{XXXL}
%  & 2000001 & XML2000000\_2121\_01 & -- & 0.2 & 204.04 & 298.89 & 5920305652.53 & \textunderscore{} \\
%  & 2000001 & XML2000000\_2123\_01 & -- & 0.2 & 172.03 & 274.20 & 1985907542.13 & \textunderscore{} \\
%  & 5000001 & XML5000000\_1173\_01 & -- & 0.1 & 458.98 & 219.36 & 2537203464.83 & \textunderscore{} \\
%  & 5000001 & XML5000000\_2176\_01 & -- & 0.1 & 492.73 & 238.55 & 1573789826.68 & \textunderscore{} \\
% \hline
% \end{tabular}
% \end{table*}

\begin{table*}[t]
\caption{Per-instance results of BP-MDS ($\rho = 10^4$).}
\centering
\small
\begin{tabular}{clrrrrrrr}
\toprule
Class & Input & N & BKS & $\alpha$ ($^\circ$) & Time (s) & Memory (MB) & Cost & Gap (\%)\\
\midrule
\multirow{2}{*}{S}
 & Golden\_17 & 241 & 707.76 & 36.0 & 0.07 & 10.52 & 762.71 & 7.76 \\
 & Golden\_18 & 301 & 995.13 & 24.0 & 0.12 & 14.88 & 1044.63 & 4.97 \\
\midrule
\multirow{5}{*}{M}
 & Leuven1 & 3001 & 192848 & 30.0 & 0.49 & 10.26 & 207573.35 & 7.64 \\
 & Leuven2 & 4001 & 111391 & 66.0 & 0.78 & 10.26 & 127544.83 & 14.50 \\
 & Antwerp1 & 6001 & 477277 & 18.0 & 0.53 & 51.85 & 512064.59 & 7.29 \\
 & Antwerp2 & 7001 & 291350 & 39.0 & 0.64 & 9.76 & 328377.84 & 12.71 \\
 & Ghent1 & 10001 & 469531 & 23.0 & 0.78 & 10.41 & 503730.90 & 7.28 \\
\midrule
\multirow{7}{*}{L}
 & Brussels1 & 15001 & 501719 & 11.5 & 1.01 & 175.10 & 544650.63 & 8.56 \\
 & Brussels2 & 16001 & 345468 & 31.0 & 2.25 & 9.75 & 390855.84 & 13.14 \\
 & Flanders1 & 20001 & 7240118 & 18.0 & 1.34 & 11.47 & 7668875.48 & 5.92 \\
 & Flanders2 & 30001 & 4373244 & 19.5 & 4.65 & 12.38 & 4864684.64 & 11.24 \\
 & XML25000\_1173\_01 & 25001 & -- & 2.0 & 2.31 & 346.74 & 46902788.87 & -- \\
 & XML50000\_2173\_01 & 50001 & -- & 2.5 & 5.04 & 343.13 & 58108045.17 & -- \\
 & Molise & 50000 & 111184982 & 23.0 & 3.35 & 19.83 & 113367183.12 & 1.96 \\
\midrule
\multirow{3}{*}{XL}
 & XML75000\_1173\_01 & 75001 & -- & 1.0 & 5.87 & 331.14 & 141858875.48 & -- \\
 & XML100000\_2173\_01 & 100001 & -- & 1.0 & 9.73 & 333.01 & 116740181.51 & -- \\
 & Trentino-Alto-Adige & 100000 & 102063181 & 34.0 & 14.52 & 34.24 & 104219633.59 & 2.11 \\
\midrule
\multirow{4}{*}{XXL}
 & Sardegna & 470000 & 827934149 & 15.0 & 132.33 & 78.73 & 828568977.54 & 0.08 \\
 & XML500000\_1173\_01 & 500001 & -- & 0.3 & 40.98 & 315.07 & 946129677.99 & -- \\
 & XML1000000\_2123\_01 & 1000001 & -- & 0.3 & 88.67 & 288.31 & 993950684.86 & -- \\
 & Lazio & 1000000 & 3145381332 & 5.0 & 101.99 & 143.92 & 3189655809.43 & 1.41\\
\midrule
\multirow{3}{*}{XXXL}
 & XML2000000\_2123\_01 & 2000001 & -- & 0.2 & 172.03 & 274.20 & 1985907542.13 & -- \\
 & XML5000000\_1176\_01 & 5000001 & -- & 0.1 & 465.79 & 218.85 & 2537175956.58 & -- \\
 & XML5000000\_2176\_01 & 5000001 & -- & 0.2 & 484.77 & 234.07 & 1573769553.19 & -- \\
\bottomrule
\end{tabular}

\label{tab:bpmds_per_instance_results}
\end{table*}


Let $k = \left\lceil \frac{360^\circ}{\alpha} \right\rceil$ denote the number of angular buckets. Let $\{s_0, s_1, \dots, \allowbreak s_{k-1}\}$ denote the sizes of the buckets, with $\sum_{j=0}^{k-1} s_j = N - 1 + k$, where the $N-1$ customers appear exactly once, while the depot is replicated across all $k$ buckets.

\textbf{Creation of Buckets:} The \textsc{Get\_Buckets} procedure assigns each of the $N-1$ customers to one of the $k$ buckets in $O((N-1)k)$ time, while bucket initialization and depot insertion require $O(k)$ time, yielding an overall time complexity of $O(Nk)$ and space complexity of $O(N + k)$.

\textbf{Per-Bucket Processing:} For a bucket of size $s_j$, MST construction (Prim’s with binary heap on an implicit complete graph) takes $O(s_j^2 \log s_j)$ time and $O(s_j)$ space. Each DFS-based route generation requires $O(s_j)$ time and space. Route refinement (e.g., 2-opt) costs $O(s_j^2)$ time and $O(s_j)$ space. 

$O\!\left(Nk + \sum_{j=0}^{k-1} \left( s_j^2 \log s_j + \rho\, s_j + s_j^2 \right)\right)$ is the overall time complexity in the sequential setting, while the space complexity is $O(N + k)$. Using the fact that $\sum_{j=0}^{k-1} s_j^2 \log s_j \le \left(\sum_{j=0}^{k-1} s_j\right)^2 \log\!\left(\sum_{j=0}^{k-1} s_j\right) = (N-1+k)^2 \log(N-1+k)$ (maximized when all the vertices lie in a single bucket), and noting that $k \ll N$, we obtain $\sum_{j=0}^{k-1} s_j^2 \log s_j = O(N^2 \log N)$. Hence, the overall sequential time complexity is upper bounded by $O(N^2 \log N)$, although this bound is conservative. In the parallel setting (assuming full parallelism), the overall time complexity is $O\!\left(Nk + \max_{0 \le j \le k-1} \left( s_j^2 \log s_j + s_j + s_j^2 \right)\right)$, where the execution time is dominated by the largest bucket.  For unbalanced buckets, the largest bucket becomes the bottleneck due to its $O(s_j^2 \log s_j)$ MST construction cost. Since the average bucket size is $\bar{s} = \frac{\sum_{j=0}^{k-1} s_j}{k} = \frac{N-1+k}{k} \approx \frac{N}{k}$, increasing $k$ reduces the average bucket size, thereby lowering the dominant per-bucket computation cost and improving parallel scalability. For approximately balanced buckets, i.e., $s_j \approx \frac{N}{k}$ for all $j$, $O\!\left(Nk + \frac{N^2}{k} \log \frac{N}{k} + \rho N  + \frac{N^2}{k} \right)$ is the sequential time complexity, while the parallel time complexity reduces to $O\!\left(Nk + \frac{N^2}{k^2} \log \frac{N}{k} + \frac{N}{k} + \frac{N^2}{k^2}\right)$. This demonstrates a near-$k$-fold reduction in the dominant quadratic term under full parallelism and eliminates the $\rho$ component. The running time decreases with increasing $k$, while the resulting increase in buckets enables greater parallelism, yielding strong scalability with increasing core counts.


% The running time decreases with increasing $k$, while the proportional increase in parallelism enables strong scalability with increasing core counts.

% The running time decreases with increasing $k$, while the degree of parallelism grows proportionally, as more buckets enable greater concurrency across threads. Consequently, the approach exhibits strong scalability with an increasing number of cores.